\begin{abstract}
\textbf{Objective:} We investigate whether
source-domain regularisation selection in
cross-subject EEG changes when an
often-used class-centred covariance
validation score is replaced by a
source-trained Gaussian predictive score.
\textbf{Methods:} We distinguish held-source
person class-recentred covariance scoring,
proper plug-in class-conditional Gaussian
predictive scoring, and held-source
discrimination. We derive the exact
nonnegative class-mean-shift penalty
separating the first two objectives,
and a finite-grid source-only certificate
for when their selected shrinkage
strength changes. Original pre-data
EEG tests are preserved; corrected
controls on Shin2017A and
BNCI2014\_002 are explicitly
retrospective. A separate public
BNCI2014\_001 protocol was
committed before its first run
with a source-only selection lock.
\textbf{Results:} Original EEGMMIDB
(\(N=20\)) showed +0.0561 mean
AUC for discrimination over the
historical class-centred criterion
(Holm \(p=0.00597\)), but five
individuals were harmed. Two
original BNCI tests were
non-significant and a prospective
filter-bank prediction failed.
The corrected predictive comparator
nearly tied discrimination on
Shin2017A (\(0.786207\) versus
\(0.785900\)) and was 0.00682 AUC
below it on BNCI2014\_002.
In the separately protocol-locked
BNCI2014\_001 experiment
(\(N=9\), 2,592 trials), discrimination,
correct predictive likelihood and
historical class-centred selection
obtained 0.789052, 0.783415
and 0.772387 AUC, respectively.
The exact source-estimand decision
breakpoint was dataset-dependent.
\textbf{Conclusion:} The class-centred
covariance score does not measure
source-mean predictive likelihood,
and can change the apparent method
ranking. Results favour clear
estimand specifications and
source-only sensitivity reporting,
not a universal optimal shrinkage
rule or deployment-level safety claim.
\end{abstract}
